\begin{beginnerstatement}[Kurz erklärt: Profilmodell]

Der \emph{Core} ist die gemeinsame Grundausstattung jedes erzeugten Projekts.
Ein \emph{Feature} ist ein einzelner technischer Baustein mit den dazugehörigen
Dateien und benötigten Abhängigkeiten. Ein \emph{Profil} ist wie eine
vorkonfigurierte Ausstattung; ein \emph{Preset} ist eine solche bereits
festgelegte Auswahl. Eine \emph{Capability} ergänzt ein Profil um eine
optionale Fähigkeit. Das \emph{Manifest} hält fest, was für ein Projekt
ausgewählt und tatsächlich aufgelöst wurde. Der \emph{Generator} erstellt das
Projekt, während der \emph{Resolver} zuvor die zulässige Kombination und ihre
Voraussetzungen ermittelt. Der daraus entstehende \emph{Scaffold} ist das
Grundgerüst aus den benötigten Ordnern und Dateien. Eine Abhängigkeit ist eine
Voraussetzung; bei einer transitiven Abhängigkeit wird eine weitere
Voraussetzung nur indirekt über einen anderen Baustein benötigt.

\end{beginnerstatement}
